% Unnumbered structural observation. No new atlas, counterexample, or movement.
\subsection*{The residual canonical order of a minimal-rank anchor is local}

Let $(A,\mathfrak m,k)$ be a complete excellent normal local surface
with algebraically closed coefficient and residue field.  Let a
prime-to-characteristic geometric Brauer class on its puncture have
index and period $M$, and let $\mathcal A_U$ be a twisted anchor of
rank exactly $M$.  Write
\[
             P_0=j_*\operatorname{End}(\mathcal A_U),
             \qquad D=P_0\otimes_A\operatorname{Frac}A.
\]
The index assertion means that $D$ is a central division algebra
of degree $M$.  Take the canonical full pullback
$P=(\pi^*P_0)^{**}$ on a regular resolution.

\paragraph{The complete order.}
The algebra $P_0$ is finite over $A$, embeds in $D$, and hence has
no idempotents other than zero and one.  We use the known
noncommutative idempotent-lifting result for a module-finite algebra
over a complete local ring: idempotents modulo the Jacobson radical
lift.  The exact source read is Leuschke and Wiegand's
\emph{Cohen--Macaulay Representations}, author draft dated 9 March
2011, Proposition 1.4, Lemma 1.6 and Corollary 1.8, printed pp.~6--9.
These statements explicitly allow noncommutative algebras.

Put $\Lambda=P_0/\mathfrak mP_0$.  It is a finite-dimensional
$k$-algebra.  Its semisimple quotient has no nontrivial idempotents:
first lift an idempotent through the nilpotent radical of $\Lambda$,
then lift it to $P_0$.  Consequently its semisimple quotient is a
division algebra.  A finite-dimensional division algebra over an
algebraically closed field is $k$: every element has a split
annihilating polynomial, and the absence of zero divisors forces it
to equal one of that polynomial's scalar roots.  We obtain
\[
             \Lambda/\operatorname{rad}\Lambda\simeq k.
\]
In particular $P_0$ is local in the noncommutative sense.  This is
not a claim that $P_0$ is commutative or that its radical vanishes.

\paragraph{An original exceptional generic point.}
Let $\eta$ be the generic point of an exceptional component, and put
$O=\mathcal O_{Y,\eta}$, with parameter $z$ and residue field $K_\eta$.
At this height-one point the full pullback is the ordinary
torsion-free pullback image.  Therefore there is a surjection
\[
 P_0\otimes_A O\longrightarrow P_\eta,
 \qquad
 \Lambda\otimes_k K_\eta\longrightarrow P_\eta/zP_\eta.
\]
The second follows because $\mathfrak m$ maps into $(z)$ and
tensoring is right exact.  The radical of $\Lambda$ is nilpotent,
so after any field extension its tensor remains nilpotent and the
quotient remains a field.  Thus $\Lambda\otimes_k K_\eta$ is local.
Its nonzero unital quotient $P_\eta/zP_\eta$ is local as well, with
semisimple quotient $K_\eta$.

This reasoning also applies after extending $K_\eta$ to the residue
spectral field and taking the image in any one multiplicity space of
a $P$-stable root anchor.  The image is a unital quotient of the
same scalar-plus-nilpotent algebra.  Its semisimple quotient is scalar;
it cannot be the full matrix algebra of size greater than one.

The height-one premise is essential.  The assertion applies to the
canonical full pullback and to strict transforms of its exceptional
generic points.  It is not asserted for a newly created exceptional
curve lying over a closed basepoint of the original evaluation map
when one retains only the ordinary pullback of the old full bundle.

\paragraph{What this says about the rank-period question.}
At a residue-order-$m_i$ component, the rank-$M$ anchor has spectral
multiplicity $M/m_i$.  If this exceeds one, minimality of the global
anchor cannot be used to deduce a full matrix residual image of the
canonical order on that space.  The preceding argument instead
forces a scalar semisimple quotient.  Every such representation has
a filtration with one-dimensional factors, but the radical action
can still constrain its invariant flags strongly, for example by
being uniserial.  No arbitrary family of invariant quotient lines,
unbounded splitting, or global period-$M$ counterexample follows
from the existence of this filtration.

Accordingly the observation does not extend the saved period-one,
rank-two overorder example to minimal rank.  It identifies what an
optimized-anchor argument would have to control: the actual radical
action and the global degrees of its invariant lattices, rather than
a residual full-matrix action inferred from rank.  It yields no new
uniform generation, cohomology, presentation, or representation bound.

\paragraph{Exact source read and status.}
The source was opened as primary author text at
\url{https://www-users.cse.umn.edu/~reiner/ReadingSeminars/McKaySeminar2016/LeuschkeWiegandBookonMcKay.pdf}.
The complete written proofs of Proposition 1.4, Lemma 1.6,
Corollary 1.8 and the adjacent Theorem 1.7 passage were freshly read.
The finite-module completeness and Wedderburn--Artin inputs remain
known recursively cited facts.  For the polynomial approximation in
the lifting proof one may retain the same initial lift throughout,
so that the approximants lie in its commutative polynomial subalgebra;
no uniqueness of arbitrary noncommutative idempotent lifts is used.
The large book was not downloaded into the project.  The observation
is an unnumbered research note, not a promoted manuscript result.
